\chapter*{Preface}
\addcontentsline{toc}{chapter}{Preface}
\markboth{Preface}{Preface}


\section*{Why Invite You to Read This Book?}

Perhaps you have never heard of Revelation. Or perhaps you know that it is the last book of the Bible, yet wonder whether its message can be trusted. You may read the Scriptures regularly but find Revelation difficult to understand.

If any of these descriptions sound like you, this book was written with you in mind.

If you have not yet encountered the Christian faith, you may assume that Revelation is a book intended only for Christians. But I would like you to know that Revelation is God's own invitation to discover the good news of His gospel. After John wrote this book, it circulated among the churches of Asia Minor. Many early Christians held fast to their faith under pressure and persecution\index{persecution}, some bearing witness to Christ even at the cost of their lives. Their faith rested on the Christ who had truly come into the world.

\section*{Why I Wrote This Book}

There are many films about end-times prophecies, as well as movies and animated productions set against the backdrop of the world's final days. Such stories naturally awaken curiosity about the world portrayed in Revelation. Some people therefore turn to the Bible to discover whether it tells us what will happen at the end of the world. A teenager once told me that this was precisely why he had begun reading the Bible: he wanted to find out whether it described the events of the world's end.

Christians, too, can become preoccupied with questions about how to interpret the dragon\index{dragon}, the beast\index{beast}, the mark, the millennium\index{millennium}, and other prophetic details. My own reading of Revelation underwent a similar turning point. I came to realize that I needed to approach the book anew, not by beginning with each individual event, but by beginning with Jesus Christ.

After all, what are the opening words of Revelation? ``The revelation of Jesus Christ.'' This is a book in which Jesus Himself speaks to us. It reveals that He reigns over history, judges evil, triumphs over Satan\index{Satan} and death, will return, and will make all things new, finally bringing the people He has redeemed\index{redemption} into eternal glory\index{glory}.

I also hope that this book will serve as a resource for Sunday school teachers. For this reason, I have sought to explain the background and message of the biblical text as clearly as possible. My prayer is that whether you are reading for personal devotion, studying Scripture in a small group, or preparing to teach others, this book will help you reflect with others on the Jesus Christ to whom Revelation bears witness.

\section*{Christ, the Lamb, at the Center: How Should We Read Revelation?}\index{Lamb, the}

As you read Revelation, you may find yourself puzzled. Wasn't Jesus a man? How, then, does He become a lamb?

To answer this question, we must learn to read the Bible as a whole.

When Jesus first came into the world, He came as the Lamb\index{Lamb, the} who takes away sin. He willingly endured the cruelest form of execution devised by human hands: crucifixion. This was no accident or helpless surrender. He gave Himself willingly to redeem those held captive by sin. How do we know that He suffered for our transgressions? Some seven hundred years before Jesus came to earth, the prophet Isaiah had already foretold His sacrificial death:

\begin{quote}
``But he was pierced for our transgressions, he was crushed for our iniquities; the punishment that brought us peace was on him, and by his wounds we are healed.''
---Isaiah 53:5 (NIV)
\end{quote}

He was led like a lamb to the slaughter, offering Himself as the sacrifice for our sins.

But did the story end there? No. Revelation carries the story of Jesus forward. He came as the Lamb who was slain\index{Lamb, the!slain}, but now He stands as the victorious Lamb, the conquering\index{overcomers} King who leads the armies of heaven. We recognize Him by His robe, dipped in blood (Revelation 19:13). The One who gave His life is the One who reigns in victory.

Revelation, then, is a story centered on Jesus Christ. He reigns, opens the scroll\index{scroll}, judges evil, defeats Satan\index{Satan}, and brings the people He has redeemed into the glory of the new creation.

How blessed are the countless Christians throughout the ages who have been brought by grace\index{grace} into the glory of God! I believe that you, too, may be among those richly blessed by His grace. My hope is that, after reading this book, you will respond to the Lord's gracious invitation and come to know the One who calls you into His glory.

\section*{About The \textit{by Grace} Series}\index{grace}

This book continues the naming theme of my earlier work, \textit{Caught by Grace}. Together, these titles trace the journey of a person saved by grace: first, caught by God's grace; then, changed by His grace; and finally, brought by His grace into His glory.

The three titles in this series express this journey:

\begin{itemize}
    \item \textit{Caught by Grace}
    \item \textit{Changed by Grace}
    \item \textit{Brought into Glory by Grace }
\end{itemize}

On the cover of this book, the word ``Grace'' appears in red, symbolizing the blood of Jesus Christ through which we are brought into His glory.

Why, then, write the third book, \textit{Brought into Glory by Grace}, first? Because \textbf{knowing the destination helps us understand the purpose of the transformation along the way.} God's salvation in Christ has a glorious end: we will ultimately share in the glory of Christ Himself.

When we know where we are going, we can better understand why God must continue to transform us today. This is the truth Paul sets forth in the golden chain of salvation in Romans 8:

\begin{quote}
``For those God foreknew he also predestined to be conformed to the image of his Son, that he might be the firstborn among many brothers and sisters. And those he predestined, he also called; those he called, he also justified; those he justified, he also glorified.''
---Romans 8:29--30 (NIV)
\end{quote}

From foreknowledge and predestination to calling, justification, and finally glorification, the journey of salvation is not a path we travel by our own strength toward glory. It is God who, by His grace, leads those He has redeemed all the way to that glorious destination.

\section*{Research Background}

The literature on Revelation is vast, and in preparing this introductory guide, I have necessarily had to be selective. In addition to consulting the work of several biblical scholars and commentators, I have drawn on my notes and teaching materials.
As I revisited and reorganized this material, I made many new discoveries and gained fresh insights. 

The following authors and works were particularly important to this book:

\begin{enumerate}
    \item John MacArthur\index{MacArthur, John}\footnote{MacArthur, J. (2007). \textit{Because the Time Is Near: John MacArthur Explains the Book of Revelation}. Moody Publishers.} --- His verse-by-verse exposition and his discussions of Christ's return, sovereign reign, and ultimate victory were valuable resources.

    \item Grant R. Osborne\index{Osborne, Grant R.}\footnote{Osborne, G. R. (2002). \textit{Revelation}. Baker Academic.} --- His examination of different approaches to interpreting Revelation helped me consider its symbols, literary structure, and eschatological themes from multiple perspectives.

    \item G. K. Beale\index{Beale, G. K.}\footnote{Beale, G. K. (1999). \textit{The Book of Revelation: A Commentary on the Greek Text}. Wm. B. Eerdmans Publishing Co.} --- His extensive research into the Old Testament background of Revelation helped me read this book within the larger biblical story and recognize the richness of its Old Testament imagery and theological connections.

    \item Craig R. Koester\index{Koester, Craig R.}\footnote{Koester, C. R. (2014). \textit{Revelation: A New Translation with Introduction and Commentary}. Yale University Press.} --- His work on the historical setting and literary context of Revelation helped me better understand the message this book conveyed to its original readers in the churches of Asia Minor.

    \item Richard Bauckham\footnote{Bauckham, R. (1993). \textit{The Theology of the Book of Revelation}. Cambridge University Press.} --- His theological studies helped me explore the book's central themes more deeply.
\end{enumerate}

These scholars make distinct contributions and offer different interpretations. I have therefore sought, as far as possible, to present the major alternative viewpoints fairly. My purpose is to guide readers through Revelation in an accessible way while keeping its central focus firmly on the victory of Jesus Christ and the ultimate entrance of His people into eternal glory.

\section*{Tables and Timelines}

Revelation contains a wealth of visions, images, and theological themes. To help modern readers grasp this richness, I have organized information into tables and timelines wherever these formats can present it clearly and concisely.

In this way, the main text can remain focused on the sovereignty and victory of Jesus Christ and His ultimate work of making all things new. Where different interpretations of a passage need further explanation, I have addressed them in footnotes or appendices.

\section*{Illustrations in This Book}

Revelation is filled with vivid and powerful images: the dragon\index{dragon}, the beasts, the Lamb, the golden lampstands\index{lampstands}, the trumpets, the seals, the bowls, and the New Jerusalem\index{New Jerusalem}. For modern readers, these visions can sometimes be difficult to picture from words alone. I have therefore included selected illustrations to help readers visualize the scenes and better grasp the message conveyed by the biblical text.

Although conceiving and producing these illustrations required considerable time and effort, I believe that effort is worthwhile if it helps brothers and sisters enter more deeply into the visionary world of Revelation.

Yet illustrations can also shape the reader's imagination, sometimes introducing details that the biblical text itself does not explicitly provide. For this reason, I have sought to keep the illustrations simple, using them primarily as visual aids rather than allowing them to replace the descriptions and explanations of Scripture.\footnote{The illustrations in this book were developed and refined by the author through repeated revisions with Google Gemini, then reviewed, edited, and selected by the author for the final manuscript.}
The illustrations are intended to assist in understanding its figures, scenes, and message, not to serve as an authoritative interpretation of Scripture.

\section*{Closing Reflections}

When I completed the first draft of this book, my heart was filled with awe. I, a sinner saved by grace, had undertaken to write, as best I could, of His glory. The radiant light of His love shone upon me, making my pen race joyfully across the page. Yet before my Creator\index{Creator}, this clumsy pen of mine, so inadequate to express what I longed to say, seemed little more than the startled cry of a crow.

But my Creator\index{Creator} knows. He knows that what I have sought to proclaim through these pages is the glory of the Lord whom I love, worship\index{worship}, and await with longing.

And so, with the apostle John, I close with these words from the final pages of Revelation:

\begin{quote}
``Amen. Come, Lord Jesus.''
---Revelation 22:20 (NIV)
\end{quote}

\vspace{2em}

\begin{flushright}
Your servant in Christ,\\
Chris Kuo\\
Elder, Reformed Church of Newtown\\
Elmhurst, New York
\end{flushright}
